\section{Preconditioning and Interface Obstructions}
\label{sec:preconditioning}

``Just precondition'' and ``just keep the state'' are incomplete complexity
arguments.  Preconditioning may exchange spectral condition for factor
normalization, right-hand-side preparation, or recovery sensitivity.  A
normalized Newton state may in turn be too weak to supply the coordinatewise
diagonal needed by the next Newton system.  This section prices those exchanges
in the access models of \cref{sec:prelim-quantum-access}.  The conclusions are
products and complete ledgers, never lower bounds on one spectral factor in
isolation.

The analytic condition/normalization statements, including the fixed-factor
minimax result, are model free.  Only their compilation/accounting
corollaries, the parity reductions, and the lazy-oracle results use the raw
sparse coefficient model of
\cref{def:raw-sparse-access}.  The state-interface results instead use the
distinct state-preparation access of \cref{def:state-preparation-access}; they
are not raw-coefficient query lower bounds.  The lazy-oracle results finally
explain why the caching ceiling in \cref{sec:lp-gain-plateau} prevents an
unconditional per-iteration multiplier on a fixed instance.

\subsection{Factor normalization and whitening}
\label{sec:preconditioning-whitening}

We first isolate the signed gain-propagation block used by
\cref{thm:lp-gain-plateau}.  Let $N\geq2$, $M=17N$, and
\[
 T=\left\lceil\tfrac12\log_2N\right\rceil,
 \qquad H=2^T,\qquad H_i=2^{\min\{i,T\}},\qquad g_i=H_i/H_{i-1},
\]
and take $a_i=\varsigma_i$ for $i\leq N$ and $a_i=1$ afterward.  Thus
$\varsigma_i$ here is the hidden sign denoted $\sigma_i$ in
\cref{thm:lp-gain-plateau}.  Here \(H\) is the plateau height; the normal
matrix \(H_\mu\) used later is unrelated.  On
$d=(d_0,\ldots,d_M)^T$, define the pinned lower-bidiagonal factor
\begin{equation}
 (B_\varsigma d)_0=d_0,
 \qquad (B_\varsigma d)_i=d_i-g_ia_id_{i-1}.
 \label{eq:prec-gain-factor}
\end{equation}
Put $\tau_0=1$, $\tau_i=\prod_{j=1}^ia_j$, and
$D_\tau=\diag(\tau_0,\ldots,\tau_M)$.  Let $B_+$ be
\eqref{eq:prec-gain-factor} with every sign positive and set
$S_+=B_+B_+^T$.  Then
\begin{equation}
 B_\varsigma=D_\tau B_+D_\tau,
 \qquad S_\varsigma:=B_\varsigma B_\varsigma^T
 =D_\tau S_+D_\tau.
 \label{eq:prec-gauge}
\end{equation}
Thus the singular values are independent of the signs, although implementing
the gauge requires their prefix products.  This switching identity and the
corresponding inverse formulas are standard signed-graph facts
\cite{alazemi2024-signed-graphs-and-inverses-of}.

\begin{lemma}[Plateau singular scale]
\label{lem:prec-small-singular}
If $L=M-T+1$ is the number of plateau nodes, then
\begin{equation}
 \sigma_{\min}(B_\varsigma)
 \leq2\sin\frac{\pi}{2(L+1)}
 \leq\frac{\pi}{L+1}<\frac{\pi}{16N},
 \qquad \norm{B_\varsigma^{-1}}_2=\Theta(N).
 \label{eq:prec-small-singular}
\end{equation}
\end{lemma}

\begin{proof}
Take a Dirichlet sine vector supported on nodes $T,\ldots,M$, with virtual
zeros at $T-1$ and $M+1$.  On this support $B_+$ is the first-difference
operator.  Adding the omitted final difference can only increase its energy,
and the full Dirichlet path quotient is
$4\sin^2(\pi/(2(L+1)))$.  Equation~\eqref{eq:prec-gauge} gives the signed
case, while $\sin x\leq x$ and $L+1>16N$ give the two remaining inequalities.

For $i\geq j$, direct forward substitution gives
\begin{equation}
 (B_+^{-1})_{ij}=H_i/H_j.                    \label{eq:prec-factor-inverse}
\end{equation}
The column $B_+^{-1}e_0=(H_i)_i$ has norm $\Theta(N)$, whereas summing the
squares in \eqref{eq:prec-factor-inverse} gives
$\norm{B_+^{-1}}_F=O(N)$.  Orthogonal equivalence in
\eqref{eq:prec-gauge} proves the final claim.
\end{proof}

\begin{theorem}[Sharp plateau whitening frontier]
\label{thm:prec-whitening-frontier}
Let $P_\varsigma$ be any square factor such that, for $K\geq1$,
\begin{equation}
 \frac1K I\preceq P_\varsigma S_\varsigma P_\varsigma^*
 \preceq I.                                      \label{eq:prec-whitened-spectrum}
\end{equation}
Then
\begin{equation}
 \norm{P_\varsigma}_2\sqrt K
 \geq\frac1{\sigma_{\min}(B_\varsigma)}
 >\frac{16N}{\pi}.                              \label{eq:prec-whitening-frontier}
\end{equation}
Consequently every exact $(\alpha_P,a_{\rm anc},0)$ block encoding of $P_\varsigma$
obeys
\begin{equation}
 \boxed{\alpha_P\sqrt K>\frac{16N}{\pi}}.       \label{eq:prec-alpha-frontier}
\end{equation}
For $1\leq a\leq\Theta(N)$, the regularized family
\begin{equation}
 P_\varsigma^{(a)}
 =D_\tau(S_++a^{-2}I)^{-1/2}D_\tau             \label{eq:prec-regularized-family}
\end{equation}
has $\norm{P_\varsigma^{(a)}}_2=\Theta(a)$ and, after a constant
input-independent normalization of the largest eigenvalue,
\begin{equation}
 K=\Theta((N/a)^2),\qquad
 \norm{P_\varsigma^{(a)}}_2\sqrt K=\Theta(N).  \label{eq:prec-tight-frontier}
\end{equation}
Thus the order of the frontier is attained throughout, not only at its
endpoints.
\end{theorem}

\begin{proof}
Equation~\eqref{eq:prec-whitened-spectrum} implies
$\sigma_{\min}(P_\varsigma B_\varsigma)\geq K^{-1/2}$.  On a right
singular vector of $B_\varsigma$ for its least singular value,
\[
 \sigma_{\min}(P_\varsigma B_\varsigma)
 \leq\norm{P_\varsigma}_2\sigma_{\min}(B_\varsigma).
\]
This and \cref{lem:prec-small-singular} prove
\eqref{eq:prec-whitening-frontier}; the normalization of an exact block
encoding is at least the encoded norm.

For tightness, the eigenvalues of the congruence in
\eqref{eq:prec-regularized-family} are
\[
 \frac{\lambda}{\lambda+a^{-2}},\qquad \lambda\in\sigma(S_+).
\]
The bidiagonal factor has bounded norm, a diagonal entry of one, and inverse
norm $\Theta(N)$ by \cref{lem:prec-small-singular}.  Hence
$\lambda_{\max}(S_+)=\Theta(1)$ and
$\lambda_{\min}(S_+)=\Theta(N^{-2})$.  It follows that the factor norm is
$\Theta(a)$ and that the ratio of the displayed largest and smallest
eigenvalues is $\Theta((N/a)^2)$ in the stated range.  This proves
\eqref{eq:prec-tight-frontier}.
\end{proof}

The theorem is a specialization, not a new general quantum-preconditioning
obstruction.  Theorem~3 and Corollary~4 of Nie, Lai, and An already prove the
separate-block-encoding limitation for arbitrary preconditioners
\cite{hantao2026-pauli-structured-preconditioning-quantum-linear}.  What is
specific here is the constant $16/\pi$, the tight interpolation
\eqref{eq:prec-regularized-family}--\eqref{eq:prec-tight-frontier}, and the
preprocessing-aware raw-query theorem below.  Directly encoding a classically
formed or algebraically regrouped product is a different access model, as both
that work and \cref{prop:no-composed-preconditioning} emphasize.

Two useful approximate operator consequences are immediate.  If the factor is
supplied through an exact $(\alpha_P,a_{\rm anc},0)$ block encoding, then, for fixed
$\delta<1$,
\begin{align}
 \norm{P_\varsigma S_\varsigma P_\varsigma^*-I}_2\leq\delta
 &\Longrightarrow
 \alpha_P\geq
 \frac{\sqrt{1-\delta}}{\sigma_{\min}(B_\varsigma)},\label{eq:prec-approx-white}\\
 \norm{P_\varsigma B_\varsigma-I}_2\leq\delta
 &\Longrightarrow
 \alpha_P\geq
 \frac{1-\delta}{\sigma_{\min}(B_\varsigma)}.\label{eq:prec-approx-inverse}
\end{align}
Indeed, the first inequality gives
$\sigma_{\min}(P_\varsigma B_\varsigma)\geq\sqrt{1-\delta}$, while the
second gives $\sigma_{\min}(P_\varsigma B_\varsigma)\geq1-\delta$; the same
singular-vector argument then applies.
Constant-error approximate whitening or inversion supplied through such an
exact factor encoding therefore has $\alpha_P=\Omega(N)$.  More generally, if
$\beta=\norm{P_\varsigma S_\varsigma P_\varsigma^*}_2$ and the congruence
has condition at most $K$, the scale-invariant statement is
\begin{equation}
 \norm{P_\varsigma}_2\sqrt{K/\beta}
 \geq\sigma_{\min}(B_\varsigma)^{-1}.          \label{eq:prec-scale-invariant}
\end{equation}
This normalization is essential: scalar rescaling defeats an unnormalized
statement.
For a left factor normalized by
$\sigma(P_\varsigma B_\varsigma)\subseteq[1/K,1]$ and supplied through an
exact $(\alpha_P,a_{\rm anc},0)$ block encoding, the same argument gives
\begin{equation}
 \alpha_PK\geq\norm{P_\varsigma}_2K>\frac{16N}{\pi}.       \label{eq:prec-left-frontier}
\end{equation}

The factor norm still does not price compilation from raw signs.  Define
\begin{equation}
 F_\varsigma=\diag(B_\varsigma^{-1},B_+^{-1}),\quad
 G_\varsigma=\diag(B_\varsigma,B_+),\quad
 \ket r=\frac{\ket{0,e_0}+\ket{1,e_0}}{\sqrt2}.            \label{eq:prec-FG}
\end{equation}
By \eqref{eq:prec-factor-inverse},
\begin{equation}
 F_\varsigma\ket r
 =\frac1{\sqrt2}\bigl((H_i\tau_i)_i\oplus(H_i)_i\bigr),
 \qquad \norm{F_\varsigma r}_2=\left(\sum_iH_i^2\right)^{1/2}=\Theta(N).
 \label{eq:prec-inverse-state}
\end{equation}
Pair the signed and unsigned coordinates and measure their side qubit in the
$X$ basis on the $16N$ plateau-copy nodes.  Its parity-signed expectation is
\begin{equation}
 \frac{16NH^2}{\sum_iH_i^2}
 >\frac{16N}{17N+4/3}>\frac9{10}.             \label{eq:prec-inverse-bias}
\end{equation}

\begin{theorem}[Reusable inverse access--normalization tradeoff]
\label{thm:prec-inverse-access}
Suppose an offline stage makes $s$ raw sign queries and produces a reusable
exact $(\alpha,a,0)$ block encoding of $\widetilde F_\varsigma$.  One call,
its inverse, or a controlled call makes at most $q$ further raw queries.
If, for fixed $\delta\leq1/20$,
\begin{equation}
 \norm{\widetilde F_\varsigma G_\varsigma-I}_2\leq\delta, \label{eq:prec-inverse-promise}
\end{equation}
then
\begin{equation}
 \boxed{s+O(q\alpha/N)=\Omega(N)},\qquad
 \boxed{Ns+q\alpha=\Omega(N^2)}.             \label{eq:prec-inverse-product}
\end{equation}
The constants hidden in $O$ and $\Omega$ are absolute.
\end{theorem}

\begin{proof}
Since $G_\varsigma F_\varsigma=I$,
\[
 \norm{(\widetilde F_\varsigma-F_\varsigma)r}_2
 \leq\delta\norm{F_\varsigma r}_2.
\]
Thus the normalized approximate-inverse state retains a constant parity bias
under the decoder in \eqref{eq:prec-inverse-bias}, and
$\norm{\widetilde F_\varsigma r}_2\geq(1-\delta)\Theta(N)$.
Fixed-point amplitude amplification prepares it with constant success using
$O(\alpha/N)$ block-encoding calls.  Including setup, the resulting parity
algorithm uses $s+O(q\alpha/N)$ raw queries.  Quantum parity needs
$\Omega(N)$ queries \cite{beals2001-quantum-lower-bounds-by-polynomials},
proving the first inequality; multiplication by $N$
gives the second.  Equation~\eqref{eq:prec-approx-inverse} also gives
$\alpha=\Omega(N)$, so the amplification count is never subconstant.
\end{proof}

Three no-go qualifications delimit these statements.  First, a
preconditioned product supplied directly can be the identity:
$P_\varsigma=B_\varsigma^{-1}$ gives $P_\varsigma B_\varsigma=I$, and exact
Schur whitening gives $P_\varsigma S_\varsigma P_\varsigma^*=I$.  No lower
bound on the product's normalization and condition alone is possible.
Second, applying $D_\tau$ is not free.  Any operator-norm-$<1/3$
implementation needs $\Omega(N)$ raw queries: apply it to
$(\ket0+\ket M)/\sqrt2$ and measure in the
$(\ket0\pm\ket M)/\sqrt2$ basis to recover endpoint parity.  Third, the
factor-independent conservation law needs the hard input/output contract.
Consider specifically the plateau solve
$G_\varsigma x=r$ from \eqref{eq:prec-FG}, and require an output state within
trace distance $1/20$ of
$F_\varsigma\ket r/\norm{F_\varsigma r}_2$.  For any such end-to-end
preconditioned implementation, compose its output with the fixed $X$-basis
decoder of \eqref{eq:prec-inverse-bias}.  The reduction gives
\begin{equation}
 q_{\rm set}+q_{\rm rhs}+q_{\rm rec}+m q_{\rm BE}=\Omega(N),
 \label{eq:prec-end-to-end-conservation}
\end{equation}
where $m$ is the number of solver calls, $q_{\rm BE}$ is the raw cost per
call, the solve phase makes no other raw queries, and the other terms charge
setup, transformed-RHS preparation, and recovery.  Indeed,
composing all phases with that decoder gives a bounded-error parity algorithm,
so their total raw-query cost is $\Omega(N)$.  Without this particular
RHS/output guarantee, no such law is claimed; a public solution may carry no
hidden information.  An orthogonal factor rotation can move parity among the
charged phases.  The statement is for the pinned plateau block, not for every
unrestricted full-KKT preconditioner.

\subsection{Congruence condition versus physical-state recovery}
\label{sec:preconditioning-congruence}

A second frontier arises even when the preconditioner is an elementary
active--inactive scaling.  Let the fixed Euclidean space split orthogonally as
$\R^r=U\oplus K$, with projectors $Q$ and $R$, and suppose
\begin{equation}
 H_\mu=\mu^{-1}C_\mu+\mu D_\mu,\qquad 0<\mu\leq\mu_0.       \label{eq:prec-congruence-H}
\end{equation}
Assume $C_\mu=QC_\mu Q$ and $C_\mu|_U\to C\succ0$;
$D_\mu\succeq0$ is uniformly bounded;
$D_{KK,\mu}\to D_{KK}\succ0$; and all other blocks converge.  For a fixed
nonzero $f\in U$, put
\begin{equation}
 p=C^{-1}f,\qquad g=D_{KK}^{-1}D_{KU}p,                    \label{eq:prec-pg}
\end{equation}
and assume the nontrivial coupling $g\ne0$.  These are fixed-instance
strict-complementarity asymptotics, not uniform claims in a growing
dimension.

For $0\leq\gamma\leq1$, define
\begin{equation}
 T_{\mu,\gamma}=\mu^{\gamma/2}Q+\mu^{-\gamma/2}R,\quad
 G_{\mu,\gamma}=T_{\mu,\gamma}H_\mu T_{\mu,\gamma},\quad
 \eta_\mu=H_\mu^{-1}f,\quad z_{\mu,\gamma}=T_{\mu,\gamma}^{-1}\eta_\mu.
 \label{eq:prec-power-congruence}
\end{equation}
Because $f\in U$, $T_{\mu,\gamma}f=\mu^{\gamma/2}f$; the normalized
transformed RHS is exactly the original RHS state.  The tradeoff below is
therefore a condition--recovery effect, not a loading effect.
The normalized contraction that recovers physical coordinates is
\begin{equation}
 \mathsf R_{\mu,\gamma}=T_{\mu,\gamma}/\norm{T_{\mu,\gamma}}_2
 =\mu^\gamma Q+R,\qquad
 \Phi_{\mu,\gamma}(w)=
 \frac{\mathsf R_{\mu,\gamma}w}{\norm{\mathsf R_{\mu,\gamma}w}_2}.
 \label{eq:prec-recovery-map}
\end{equation}
For a unit $w$, measure its local normalized-state sensitivity by
\begin{equation}
 \chi_{\mu,\gamma}(w)=
 \left\|D\Phi_{\mu,\gamma}(w)|_{w^\perp}\right\|_{2\to2}.
 \label{eq:prec-recovery-sensitivity}
\end{equation}

\begin{theorem}[Power congruence frontier]
\label{thm:prec-congruence-frontier}
Under the preceding assumptions, for each fixed $\gamma\in[0,1]$ as
$\mu\downarrow0$,
\begin{align}
 \kappa_2(G_{\mu,\gamma})
 &=\Theta(\mu^{-2(1-\gamma)}),                         \label{eq:prec-congruence-kappa}\\
 \eta_\mu&=\mu\binom{p}{-g}+o(\mu),                    \label{eq:prec-physical-expansion}\\
 z_{\mu,\gamma}&=
 \binom{\mu^{1-\gamma/2}(p+o(1))}
       {-\mu^{1+\gamma/2}(g+o(1))}.                    \label{eq:prec-scaled-expansion}
\end{align}
Writing $\widehat z=z/\norm z_2$, exact normalized recovery holds and direct
postselection has
\begin{equation}
 \Phi_{\mu,\gamma}(\widehat z)
 =\frac{\eta_\mu}{\norm{\eta_\mu}_2},\qquad
 P_{\rm rec}=\norm{\mathsf R_{\mu,\gamma}\widehat z}_2^2
 =\Theta(\mu^{2\gamma}),                                \label{eq:prec-recovery-probability}
\end{equation}
while
\begin{equation}
 \chi_{\mu,\gamma}(\widehat z)=\Theta(\mu^{-\gamma}).   \label{eq:prec-recovery-chi}
\end{equation}
Consequently
\begin{equation}
 \boxed{\sqrt{\kappa_2(G_{\mu,\gamma})}\,
 \chi_{\mu,\gamma}(\widehat z)=\Theta(1/\mu)},
 \qquad
 \sqrt{\kappa_2(G_{\mu,\gamma})}P_{\rm rec}^{-1/2}
 =\Theta(1/\mu).                                         \label{eq:prec-congruence-product}
\end{equation}
\end{theorem}

\begin{proof}
Set $a_\mu=\mu^{-(1-\gamma)}$ and
$b_\mu=\mu^{1-\gamma}=a_\mu^{-1}$.  In $U\oplus K$ blocks,
\begin{equation}
 G_{\mu,\gamma}=
 \begin{pmatrix}
 a_\mu C_\mu+\mu^{1+\gamma}D_{UU,\mu}&\mu D_{UK,\mu}\\
 \mu D_{KU,\mu}&b_\mu D_{KK,\mu}
 \end{pmatrix}.                                         \label{eq:prec-congruence-blocks}
\end{equation}
The upper-left block $A_\mu$ is $\Theta(a_\mu)I_U$, and its Schur
complement is
\[
 b_\mu D_{KK,\mu}-\mu^2D_{KU,\mu}A_\mu^{-1}D_{UK,\mu}
 =\Theta(b_\mu)I_K.
\]
The elimination factor has off-diagonal norm
$O(\mu b_\mu)=o(1)$.  A uniformly conditioned near-identity congruence
therefore reduces \eqref{eq:prec-congruence-blocks} to blocks of sizes
$\Theta(a_\mu)$ and $\Theta(b_\mu)$, proving
\eqref{eq:prec-congruence-kappa}.

The $K$ block of $H_\mu\eta_\mu=f$ gives exactly
\[
 \eta_{K,\mu}=-D_{KK,\mu}^{-1}D_{KU,\mu}\eta_{U,\mu}.
\]
Substitution in the $U$ block gives
\[
 [\mu^{-1}C_\mu+\mu E_\mu]\eta_{U,\mu}=f,
 \quad
 E_\mu=D_{UU,\mu}-D_{UK,\mu}D_{KK,\mu}^{-1}D_{KU,\mu},
\]
where $E_\mu$ is bounded.  Hence
$\eta_{U,\mu}=\mu(C_\mu+\mu^2E_\mu)^{-1}f=\mu(p+o(1))$,
which proves \eqref{eq:prec-physical-expansion}; applying $T^{-1}$ proves
\eqref{eq:prec-scaled-expansion}.

Because $\mathsf Rz=\mu^{\gamma/2}\eta$, normalized recovery is exact.
The two expansions give
\[
 \norm{\mathsf R\widehat z}_2
 =\mu^\gamma
 \frac{(\norm p_2^2+\norm g_2^2+o(1))^{1/2}}
      {(\norm p_2^2+\mu^{2\gamma}\norm g_2^2+o(1))^{1/2}}
 =\Theta(\mu^\gamma),
\]
proving \eqref{eq:prec-recovery-probability}.  Differentiation gives, with
$\widehat\eta=\Phi(w)$,
\begin{equation}
 \chi_{\mu,\gamma}(w)=
 \frac{\norm{(I-\widehat\eta\widehat\eta^T)
 \mathsf R_{\mu,\gamma}(I-ww^T)}_2}
      {\norm{\mathsf R_{\mu,\gamma}w}_2}.                \label{eq:prec-chi-formula}
\end{equation}
This is $O(\mu^{-\gamma})$ at $w=\widehat z$.  For $\gamma>0$, project
$g/\norm g$ onto $\widehat z^\perp$.  The resulting unit tangent is
$g/\norm g+o(1)$ in $K$, while the recovered state tends to
$(p,-g)/(\norm p^2+\norm g^2)^{1/2}$.  Since $p\ne0$, the numerator in
\eqref{eq:prec-chi-formula} is bounded below by
$\norm p/(\norm p^2+\norm g^2)^{1/2}>0$.  This proves the matching lower
bound.  For $\gamma=0$, $\mathsf R=I$ and $\chi=1$.  Multiplying the
estimates proves \eqref{eq:prec-congruence-product}.
\end{proof}

The same two-sparse LP witnesses every exponent.  Consider
\begin{equation}
 \min_{x\geq0}x_2+x_3,\qquad
 x_1+x_2=1,\quad x_2-x_3=0,                               \label{eq:prec-small-witness}
\end{equation}
so
$A=\left(\begin{smallmatrix}1&1&0\\0&1&-1\end{smallmatrix}\right)$,
$b=e_1$, and $c=(0,1,1)^T$.  Its unique optimum is $(1,0,0)$.  Along the
central path $x_2=x_3=t_\mu$, $x_1=1-t_\mu$, with
$t_\mu=\mu+O(\mu^2)$, and
\begin{equation}
 H_\mu=A(XS^{-1})A^T
 =\frac{(1-t_\mu)^2}{\mu}
 \begin{pmatrix}1&0\\0&0\end{pmatrix}
 +\frac{t_\mu^2}{\mu}
 \begin{pmatrix}1&1\\1&2\end{pmatrix}.                 \label{eq:prec-witness-normal}
\end{equation}
Indeed, dual feasibility and complementarity give equal inactive slacks
$q_\mu=1+\mu/[2(1-t_\mu)]$ and $t_\mu q_\mu=\mu$, which yields the
expansion.  For $U=\operatorname{span}\{e_1\}$ and
$K=\operatorname{span}\{e_2\}$, the limiting blocks in
\eqref{eq:prec-congruence-H} are $C=e_1e_1^T$ and
$D=\left(\begin{smallmatrix}1&1\\1&2\end{smallmatrix}\right)$.  Thus
$p=1$, $g=1/2$, and $H_\mu^{-1}b=\mu(1,-1/2)^T+o(\mu)$.
This fixed two-row, three-variable LP therefore realizes all parts of
\cref{thm:prec-congruence-frontier}; at $\gamma=1$ its scaled condition has
a finite positive limit and its recovery amplitude is asymptotic to
$\sqrt{5/4}\,\mu$.

The power family is also the complete order-optimal interpolation within all
scalar block-diagonal SPD congruences.  Let
\begin{equation}
 T_\mu=t_U(\mu)Q+t_K(\mu)R,\qquad
 r_\mu=t_U/t_K,\quad
 M_\mu=\max\{r_\mu/\mu,\mu/r_\mu\},\quad
 K_\mu=\max\{r_\mu,r_\mu^{-1}\}.                         \label{eq:prec-general-scalars}
\end{equation}
Put
\begin{equation}
 \widetilde G_\mu=T_\mu H_\mu T_\mu,\qquad
 \widetilde z_\mu=T_\mu^{-1}\eta_\mu,\qquad
 \widetilde{\mathsf R}_\mu=T_\mu/\norm{T_\mu}_2,\qquad
 \widetilde\Phi_\mu(w)=
 \frac{\widetilde{\mathsf R}_\mu w}
      {\norm{\widetilde{\mathsf R}_\mu w}_2}.            \label{eq:prec-general-recovery}
\end{equation}
Define $\widetilde P_{\rm rec}$ by applying the squared norm in
\eqref{eq:prec-recovery-probability} to
$\widetilde{\mathsf R}_\mu$ and $\widehat{\widetilde z}_\mu$, and define
$\widetilde\chi_\mu$ by \eqref{eq:prec-recovery-sensitivity} with
$\widetilde\Phi_\mu$.

\begin{theorem}[General scalar-block frontier]
\label{thm:prec-general-scalar-frontier}
Uniformly over all positive scalar functions $t_U,t_K$,
\begin{equation}
 \sqrt{\kappa_2(\widetilde G_\mu)}=\Theta(M_\mu),\quad
 \widetilde P_{\rm rec}^{-1/2}=\Theta(K_\mu),\quad
 \widetilde\chi_\mu=\Theta(K_\mu).                      \label{eq:prec-general-frontier}
\end{equation}
Consequently
\begin{equation}
 \sqrt{\kappa_2(\widetilde G_\mu)}\,\widetilde\chi_\mu
 =\Theta(M_\mu K_\mu)\geq c/\mu.                       \label{eq:prec-general-product}
\end{equation}
The same holds with $\widetilde P_{\rm rec}^{-1/2}$ in place of
$\widetilde\chi_\mu$.  Equality in order,
$M_\mu K_\mu=\Theta(1/\mu)$, holds exactly for the maximal class
\begin{equation}
 r_\mu=\Omega(\mu),\qquad r_\mu=O(1).                    \label{eq:prec-equality-class}
\end{equation}
More explicitly,
\begin{equation}
 M_\mu K_\mu=
 \begin{cases}
  \mu/r_\mu^2,&r_\mu<\mu,\\
  1/\mu,&\mu\leq r_\mu\leq1,\\
  r_\mu^2/\mu,&r_\mu>1.
 \end{cases}                                             \label{eq:prec-three-cases}
\end{equation}
\end{theorem}

\begin{proof}
The assumptions imply uniform spectral equivalence
\begin{equation}
 c_0(\mu^{-1}Q+\mu R)\preceq H_\mu
 \preceq C_0(\mu^{-1}Q+\mu R).                           \label{eq:prec-spectral-equivalence}
\end{equation}
For the lower bound, write a vector as $u+k$.  Fixed $c,d,M>0$ give
\[
 (u+k)^TH_\mu(u+k)
 \geq(c/\mu-2M^2\mu/d)\norm u^2+(d\mu/2)\norm k^2,
\]
after Young's inequality; for small $\mu$ this proves the lower half of
\eqref{eq:prec-spectral-equivalence}, and boundedness proves the upper half.
Congruence by $T_\mu$ shows that the two spectral scales are
$t_U^2/\mu$ and $\mu t_K^2$, proving the condition estimate in
\eqref{eq:prec-general-frontier}.

Equation~\eqref{eq:prec-physical-expansion} gives
$\norm{\eta_{U,\mu}}=\Theta(\mu)$ and
$\norm{\eta_{K,\mu}}=\Theta(\mu)$.  Therefore
\[
 \norm{\widetilde{\mathsf R}_\mu\widehat{\widetilde z}_\mu}
 =\frac{\norm{\eta_\mu}}
 {\max\{t_U,t_K\}
  \sqrt{\norm{\eta_{U,\mu}}^2/t_U^2+
        \norm{\eta_{K,\mu}}^2/t_K^2}}
 =\Theta(K_\mu^{-1}).
\]
This proves the recovery-probability claim.  In the plane spanned by the
normalized $U$ and $K$ components, differentiating the output angle of
$\diag(t_U,t_K)$ gives tangent gain
\[
 \frac{r_\mu^{-1}\norm{\eta_{U,\mu}}^2+
       r_\mu\norm{\eta_{K,\mu}}^2}
      {\norm{\eta_{U,\mu}}^2+\norm{\eta_{K,\mu}}^2}
 =\Theta(K_\mu).
\]
This is a lower bound on $\widetilde\chi_\mu$; differentiation as in
\eqref{eq:prec-chi-formula} gives the matching upper bound.  The elementary
three-case calculation \eqref{eq:prec-three-cases} proves the product and
the equality class.
\end{proof}

This is not a universal preconditioning lower bound.  It excludes arbitrary
left preconditioners, input-dependent changes mixing $U$ and $K$, direct
preparation of the physical state, variable-time recovery, and algorithms
that request only an observable.  It also charges neither construction of
$Q$ nor block-encoding normalization.  If $g=0$, the recovery conclusion can
fail.  In particular, the state-preserving left equilibration in
\cref{thm:stable-split-equilibration} lies outside this congruence-recovery
contract and is a genuine positive escape.  The theorem is therefore a sharp
structural lemma for a specified family, realized by
\eqref{eq:prec-small-witness}, rather than a claim against all
preconditioners.

\subsection{The parity preconditioner dichotomy}
\label{sec:preconditioning-parity}

The parity LP gives a complementary dichotomy.  Its reduced primal Hessian is
condition one, but its equality normal matrix is a signed grounded-tree
Laplacian.  Fixed preconditioners cannot serve all signings, whereas a perfect
data-dependent sparse factor is locally trivial.  The global information then
appears in factor application, transformed-RHS preparation, or recovery.

The minimax claim already holds on a path.  Let
\begin{equation}
 B=\begin{pmatrix}
 1\\-1&1\\&-1&1\\&&\ddots&\ddots
 \end{pmatrix},\qquad C=BB^T,\qquad
 R=\diag(1,-1,1,-1,\ldots)                               \label{eq:prec-path-pair}
\end{equation}
for $m\geq1$.  Both $C$ and $RCR$ are legal sign inputs.

\begin{theorem}[Minimax fixed-preconditioner bound]
\label{thm:prec-fixed-minimax}
For SPD matrices write $\kappa=\lambda_{\max}/\lambda_{\min}$.
For every input-independent SPD matrix $M$,
\begin{equation}
 \max\left\{
 \kappa(M^{-1/2}CM^{-1/2}),
 \kappa(M^{-1/2}RCRM^{-1/2})
 \right\}\geq\frac{4m^2-1}{3}\geq m^2.                  \label{eq:prec-minimax}
\end{equation}
The identity gives condition number at most $4m^2$ for either signing, so it
is minimax optimal in order.  The statement allows dense $M$ and holds more
generally for every fixed invertible factor $P$, with the two matrices
$P^TCP$ and $P^TRCRP$.  The bound $(4m^2-1)/3$ is a witness lower bound,
not a claim of the exact minimax value.  It strengthens the former
$m^2/4$ bound for $m\geq4$.
\end{theorem}

\begin{proof}
For SPD $A_0,A_1,M$, generalized Rayleigh quotients give
\begin{equation}
 \kappa(A_0^{-1/2}A_1A_0^{-1/2})
 \leq\kappa(M^{-1/2}A_0M^{-1/2})
       \kappa(M^{-1/2}A_1M^{-1/2}).                       \label{eq:prec-simultaneous}
\end{equation}
Indeed, bound the numerator and denominator of
$x^TA_1x/(x^TA_0x)$ by the extremal eigenvalues in $M$ coordinates; the
independent scales cancel in the quotient of extrema.

For $r=(m,m-1,\ldots,1)^T$,
\[
 x^TCx=\sum_{j=0}^{m-2}(x_j-x_{j+1})^2+x_{m-1}^2,
 \qquad r^TCr=m.
\]
Summing the alternating witness exactly gives
\[
 r^TRCRr=\sum_{k=1}^m(2k-1)^2=\frac{m(4m^2-1)}3.
\]
Put $t=(4m^2-1)/3$.  The generalized quotient is $t$ at $r$ and
$1/t$ at $Rr$, so
$\kappa(C^{-1/2}RCRC^{-1/2})\geq t^2$.
Apply \eqref{eq:prec-simultaneous} with $A_0=C$ and $A_1=RCR$; one of its
two factors must be at least $t$.  The same argument in coordinates
$P^{-1}r$ and $P^{-1}Rr$ proves the assertion for any invertible $P$,
without a polar-decomposition assumption.

For the upper bound, write $d_i=x_i-x_{i+1}$ for $i<m-1$ and
$d_{m-1}=x_{m-1}$.  Then $x_i=\sum_{j=i}^{m-1}d_j$ and
$x^TCx=\sum_jd_j^2$.  Cauchy--Schwarz gives
$\norm{x}_2^2\leq m^2x^TCx$, while
$(a-b)^2\leq2a^2+2b^2$ gives $x^TCx\leq4\norm{x}_2^2$.
Thus $\kappa(C)\leq4m^2$.  Since $R$ is orthogonal, the same bound
holds for $RCR$; more generally every signed path is an orthogonal
diagonal conjugate of $C$.
\end{proof}

The fixed-preconditioner result is formalized in Lean, including the
actual incidence matrices, their positive definiteness, extremal-eigenvalue
condition numbers, the simultaneous-preconditioning inequality, the exact
witness ratio, and the identity upper bound for every edge signing.
The development also proves the arbitrary-factor statement directly and
identifies the SPD specialization with the positive inverse square root.
It does not formalize the data-dependent factor or quantum-query results
that follow.  The repository's \texttt{formal/PRECONDITIONER.md} maps these
claims to the proofs and records the verification commands.

For data-dependent preconditioning, let $D_\varsigma$ be the square grounded
signed incidence matrix of the chain-plus-copy tree: its root row is $e_0^T$
and the row of a nonroot node $j$ is
$e_j^T-\tau_je_{\pi(j)}^T$, where $\tau_j$ is its edge sign.  With pair
variables the equality matrix is
\begin{equation}
 A_\varsigma=D_\varsigma[I,-I],\qquad
 C_\varsigma=D_\varsigma D_\varsigma^T.                  \label{eq:prec-tree-normal}
\end{equation}
Every row of $D_\varsigma$ has at most two nonzeros, each determined by one
local sign, yet
\begin{equation}
 D_\varsigma^{-1}C_\varsigma D_\varsigma^{-T}=I,\qquad
 D_\varsigma^{-1}e_0=(p_j)_j,                             \label{eq:prec-perfect-local}
\end{equation}
where $p_j$ is the root-to-$j$ sign product.  Set
$R_\varsigma=\diag(p_j)$, and let $D_+$ be the unsigned grounded incidence
matrix on the same rooted tree.  The switching identity is
$D_\varsigma=R_\varsigma D_+R_\varsigma$.  Thus construction of a sparse
perfect preconditioner is not hard; applying its inverse exposes all prefixes.
This distinction is the classical requirement that a useful preconditioner be
cheap to apply, expressed here in support-theory language
\cite{boman2003-support-theory-for-preconditioning}.

The right-hand side $e_0$ in \eqref{eq:prec-perfect-local} occurs in an actual
Newton equation.  At the infeasible-start point
$x=s=\mathbf1$, $y=0$, $\mu=1$ for the pair LP with $c=\mathbf1$ and
$b=e_0$, the point is centered and dual feasible, and $A_\varsigma\mathbf1=0$.
With the residual convention of \eqref{eq:newton-residuals}, any centering
target $0<\sigma_c\leq1$ gives
\begin{equation}
 A_\varsigma\Delta x=e_0,\quad
 A_\varsigma^T\Delta y+\Delta s=0,\quad
 \Delta x+\Delta s=(\sigma_c-1)\mathbf1,\quad
 A_\varsigma A_\varsigma^T\Delta y=e_0.                  \label{eq:prec-actual-newton}
\end{equation}
Since $A_\varsigma A_\varsigma^T=2D_\varsigma D_\varsigma^T$, the signed
incidence congruence transforms this exact Newton RHS into
$D_\varsigma^{-1}e_0$ up to a public scalar.

For the state lower bound, take $N\geq2$ and attach two isomorphic complete
copy trees, one to the root and one to the endpoint of the $N$-edge hidden
spine.  Each has $M$ leaves and $J=2M-2$ new nodes, where
\begin{equation}
 M=2^{\lceil\log_2(16(N+1))\rceil}.                       \label{eq:prec-two-tree-size}
\end{equation}
Here $M$ and $m=N+1$ below are local to this amplified-tree construction and
are unrelated to the plateau length $M$ and path dimension $m$ used earlier.
All new edges have positive sign; write $\widehat D_\varsigma$ for this
grounded incidence factor and define
\begin{equation}
 \widehat A_\varsigma=\widehat D_\varsigma[I,-I],\qquad
 \widehat C_\varsigma=\widehat D_\varsigma\widehat D_\varsigma^T.
 \label{eq:prec-amplified-normal}
\end{equation}
At the same infeasible-start point $x=s=\mathbf1$, $y=0$, $\mu=1$, now for
the amplified pair LP with $c=\mathbf1$ and $b=e_0$, the Newton equations are
\begin{equation}
 \widehat A_\varsigma\Delta x=e_0,\quad
 \widehat A_\varsigma^T\Delta y+\Delta s=0,\quad
 \Delta x+\Delta s=(\sigma_c-1)\mathbf1,\quad
 \widehat A_\varsigma\widehat A_\varsigma^T\Delta y=e_0.
 \label{eq:prec-amplified-newton}
\end{equation}
Thus the signed-incidence congruence maps this actual amplified Newton RHS to
$\widehat D_\varsigma^{-1}e_0$ up to a public scalar.

The weighted factor below is also the exact local factor away from the start.
At any positive interior point let
$\Theta=XS^{-1}=\diag(\Theta_+,\Theta_-)$, where
$\Theta_\pm=\diag(\theta_j^\pm)_j$, and set
\begin{equation}
 \Omega=\diag(\omega_j)_j,\qquad
 \omega_j=\theta^+_j+\theta^-_j>0.                       \label{eq:prec-local-omega}
\end{equation}
Then
\begin{equation}
 \widehat A_\varsigma\Theta\widehat A_\varsigma^T
 =\widehat D_\varsigma\Omega\widehat D_\varsigma^T
 =L_\varsigma L_\varsigma^T,\qquad
 L_\varsigma:=\widehat D_\varsigma\Omega^{1/2}.         \label{eq:prec-local-factor}
\end{equation}
Here $g$ below is the corresponding eliminated Newton RHS.  In particular,
$4/5\leq x_i,s_i\leq6/5$ implies
$4/3\leq\omega_j\leq3$, hence
$\Gamma:=\max_j\omega_j/\min_j\omega_j\leq9/4<4$.

\begin{theorem}[Factor-specific Newton-RHS hardness]
\label{thm:prec-factor-rhs-hardness}
Every unconditional algorithm that, on each sign input, outputs a state within
trace distance $1/20$ of
\begin{equation}
 \ket{z_\varsigma}=
 \frac{\widehat D_\varsigma^{-1}e_0}
      {\norm{\widehat D_\varsigma^{-1}e_0}_2}             \label{eq:prec-prefix-state}
\end{equation}
makes $\Omega(N)$ raw coefficient queries.

The conclusion remains true for the weighted nearby factor $L_\varsigma$ in
\eqref{eq:prec-local-factor} and eliminated RHS $g$ whenever
\begin{equation}
 \frac{\max_j\omega_j}{\min_j\omega_j}\leq4,\qquad
 \norm{g-e_0}_1\leq\frac1{100};                           \label{eq:prec-l1-robust}
\end{equation}
the target is the normalized state proportional to
$\Omega^{-1/2}\widehat D_\varsigma^{-1}g$.  If $g$ or $\Omega$ is
input-dependent, all raw queries used to construct and access it are included.
Thus the theorem applies at any nearby central iterate satisfying these two
promises.  Freely supplied side oracles whose numerical values themselves
encode parity are outside the theorem.
\end{theorem}

\begin{proof}
In the exact case, the prefix vector is $+1$ on every node of the root copy
tree and endpoint parity $p_N$ on every node of the endpoint copy tree.  Pair
corresponding nodes and measure their side qubit in the $X$ basis.  If
$m=N+1$, then $J\geq31m$ and the probability of a paired node is
\begin{equation}
 \frac{2J}{m+2J}>\frac{62}{63}.                           \label{eq:prec-copy-mass}
\end{equation}
Conditional on that event the measurement returns $p_N$ exactly; outputting a
fair bit otherwise succeeds with probability greater than
$1/2+31/63=125/126$.  Trace error $1/20$ leaves success above $2/3$, so the
fixed decoder computes parity and requires $\Omega(N)$ queries.

For \eqref{eq:prec-l1-robust}, first take $g=e_0$ and put
$\Gamma=\max\omega/\min\omega\leq4$.  Paired nodes have squared mass at
least
\[
 \alpha\geq\frac{2J}{2J+\Gamma m}\geq\frac{31}{33}.
\]
The amplitudes of a corresponding pair are $a$ and $p_Nb$ with
$a/b\in[1/2,2]$.  The $X$ measurement is correct with probability
$1/2+ab/(a^2+b^2)\geq9/10$.  The full ideal success is therefore at least
$1/2+(2/5)(31/33)=289/330$.

For $\delta=g-e_0$, every coordinate of
$\widehat D_\varsigma^{-1}\delta$ is a signed ancestor sum, so, if $P$ is
the number of tree nodes,
\[
 \norm{\widehat D_\varsigma^{-1}\delta}_2
 \leq\sqrt P\norm\delta_1,
 \qquad
 \frac{\norm{\Omega^{-1/2}\widehat D_\varsigma^{-1}\delta}_2}
 {\norm{\Omega^{-1/2}\widehat D_\varsigma^{-1}e_0}_2}
 \leq\sqrt\Gamma\norm\delta_1\leq\frac1{50}.
\]
The corresponding normalized pure states are at Euclidean distance at most
$1/25$.  Including trace error $1/20$ changes outcome probabilities by at
most $9/100$; subtracting this from $289/330$ leaves success above $2/3$.
Parity again proves the claim.
\end{proof}

The $\ell_1$ hypothesis is essential to this proof.  On the signed spine set
$\widetilde d_j=p_j(1-j/N)$, keep the root copy tree at value one, and set the
endpoint copy tree to zero.  Then
\begin{equation}
 \norm{\widehat D_\varsigma\widetilde d-e_0}_2=N^{-1/2},  \label{eq:prec-l2-failure}
\end{equation}
although the amplified endpoint has vanished and $\widetilde d$ stays a
constant relative distance from the prefix vector.  Thus even a vanishing
$\ell_2$ feasibility residual does not imply the state closeness used above;
\eqref{eq:prec-l2-failure} is a failure of this proof route, not a no-go for
every possible $\ell_2$-robust reduction.

Nor is transformed-RHS hardness invariant under the factor.  On the amplified
tree, let $\widehat R_\varsigma=\diag(p_j)$ be the prefix gauge and let
$\widehat D_+$ be the unsigned grounded incidence matrix.  Then
$\widehat D_\varsigma=\widehat R_\varsigma\widehat D_+
\widehat R_\varsigma$, and the equally perfect factor
\begin{equation}
 \widehat P_\varsigma=\widehat D_+^{-1}\widehat R_\varsigma
 =\widehat R_\varsigma\widehat D_\varsigma^{-1}
 \label{eq:prec-orthogonal-loophole}
\end{equation}
satisfies
\begin{equation}
 \widehat P_\varsigma\widehat C_\varsigma\widehat P_\varsigma^T=I,
 \qquad
 \widehat P_\varsigma e_0=\widehat D_+^{-1}e_0=\mathbf1.
 \label{eq:prec-easy-rhs}
\end{equation}
The RHS is now input independent.  Recovery uses
$\widehat P_\varsigma^T=\widehat R_\varsigma\widehat D_+^{-T}$, so the
parity cost has moved to factor application or original-coordinate recovery.
This orthogonal-factor
loophole is why \cref{thm:prec-factor-rhs-hardness} is deliberately
factor-specific.

The invariant statement is end to end.  Specialize
\eqref{eq:prec-amplified-newton} to $\sigma_c=1$ and take the full step.  With
$p=\widehat D_\varsigma^{-1}e_0$, the unique minimum-norm direction and its
update are
\begin{equation}
 \Delta x=\tfrac12(p,-p),\qquad \Delta s=-\Delta x,\qquad
 x^+=\mathbf1+\Delta x.                                  \label{eq:prec-hard-update}
\end{equation}
Every pair of $x^+$ is $(3/2,1/2)$ or its swap, and
$X^+s^+=(3/4)\mathbf1$; this is the exact feasible center at $\mu=3/4$.
The endpoint copy tree occupies more than $15/32$ of all pairs.  Conditional
on measuring such a pair, the computational-basis coordinate identifies
$p_N$ with probability $9/10$.  Returning a fair bit otherwise gives ideal
success greater than $1/2+(2/5)(15/32)=11/16$, and trace error $1/100$
leaves success above $2/3$.

\begin{theorem}[Preconditioner-oracle accounting]
\label{thm:prec-parity-ledger}
Let $P_\varsigma$ be any possibly data-dependent invertible factor for the
normal equation in \eqref{eq:prec-amplified-newton}, with no sparsity or locality
assumption, and let an end-to-end Newton procedure output a state within trace
distance $1/100$ of the normalized $x^+$ in
\eqref{eq:prec-hard-update}.  Partition all raw coefficient queries into
setup, transformed-RHS preparation, iterative solve, and original-coordinate
recovery/update.  Then
\begin{equation}
 \boxed{Q_{\rm set}+Q_{\rm rhs}+Q_{\rm solve}+Q_{\rm rec}=\Omega(N).}
 \label{eq:prec-parity-ledger}
\end{equation}
If the preconditioned system has condition at most $K$ and the solve phase's
only raw queries occur inside at most $T(K,\epsilon)$ calls to its
preconditioned oracles, each costing at most $q_{\rm on}$ raw queries, then
\begin{equation}
 Q_{\rm set}+Q_{\rm rhs}+T(K,\epsilon)q_{\rm on}+Q_{\rm rec}=\Omega(N).
 \label{eq:prec-condition-ledger}
\end{equation}
In particular, for
$T(K,\epsilon)=\widetilde O(K\log(1/\epsilon))$ and hidden polylogarithmic
factor $L(N,\epsilon)$,
\begin{equation}
 Q+O(Kq_{\rm on}L(N,\epsilon))=\Omega(N),\qquad
 Q:=Q_{\rm set}+Q_{\rm rhs}+Q_{\rm rec}.                 \label{eq:prec-condition-corollary}
\end{equation}
\end{theorem}

\begin{proof}
All four phases together form one coefficient-query algorithm.  Compose its
output with the fixed decoder following \eqref{eq:prec-hard-update}; this
computes parity of the $N$ hidden spine signs with bounded error.  Each LP
coefficient query is simulated by $O(1)$ sign queries, proving
\eqref{eq:prec-parity-ledger}.  Under the oracle-only solve hypothesis,
$Q_{\rm solve}\leq T(K,\epsilon)q_{\rm on}$, which gives
\eqref{eq:prec-condition-ledger}, and substitution of the stated solver call
bound gives \eqref{eq:prec-condition-corollary}.
\end{proof}

A hybrid solver that also queries raw data directly remains covered by
\eqref{eq:prec-parity-ledger}, where such queries stay in $Q_{\rm solve}$.

The ledger is tight and is not a quantum--classical separation.  A classical
algorithm reads the $N$ signs, computes every prefix, and stores the gauge in
$\Theta(N)$ queries and work.  Ladner--Fischer prefix scan reduces parallel
depth to $O(\log N)$ with $O(N)$ bounded-fan-in work
\cite{ladner1980-parallel-prefix-computation}.  Caching then removes further
hidden-sign queries, but it does not make the unsigned triangular operations
free: applying $\widehat D_+^{-1}$ or $\widehat D_+^{-T}$ still takes
$\Theta(N)$ work and $O(\log N)$ parallel depth by prefix or subtree
aggregation.  Multiplication by the cached diagonal gauge is local.  Quantum
parity changes only the constant in the query count.  The right referee-level
interpretation is therefore a
lemma/counterexample suite: \cref{thm:prec-fixed-minimax} is a finite-dimensional
minimax lower bound, \cref{thm:prec-factor-rhs-hardness} locates cost for a canonical factor,
and \cref{thm:prec-parity-ledger} records how another factor can relocate it.
It is not a standalone universal preconditioning lower bound.  Contemporary
preconditioned QIPM analyses that assume prepared transformed data are
consistent with this conclusion; reducing those oracles to raw LP access is a
separate charged step \cite{zeguan2026-preconditioned-inexact-infeasible-quantum-interi}.

\subsection{Newton-state output is not an iterate oracle}
\label{sec:preconditioning-state-interface}

We now change oracle models.  Let $v\in\R_{>0}^n$, let its known norm be
$r_v=\norm v_2$, and grant controlled state-preparation access
\begin{equation}
 U_v\ket0=\ket{\widehat v}=\frac1{r_v}\sum_{j=1}^nv_j\ket j,
 \label{eq:prec-state-oracle}
\end{equation}
together with $U_v^\dagger$ and a fixed real phase.  This is precisely
\cref{def:state-preparation-access}; it is stronger than receiving copies but
distinct from coordinate-value, raw-coefficient, QRAM, or already compiled
diagonal access.  Those interfaces are not generally ordered by strength.
The next results count calls to $U_v,U_v^\dagger$, not raw LP queries.

\begin{theorem}[State-to-diagonal normalization--query product]
\label{thm:prec-state-diagonal}
Suppose a controlled circuit $W_v$, using $q$ calls to
$U_v,U_v^\dagger$, is an $(\alpha,a,1/8)$ block encoding of
$\diag(v)$.  If the reported normalization satisfies the uniform known bound
$\alpha\leq A$, where $A\geq1$, then in the worst case
\begin{equation}
 \boxed{qA=\Omega(\sqrt n)}.                              \label{eq:prec-qA}
\end{equation}
The product is tight in order on the hard near-uniform family.
\end{theorem}

\begin{proof}
For $n\geq10$, take the positive equal-norm pair
\begin{equation}
 v^{(0)}=(1,\ldots,1),\qquad
 v^{(1)}=(2,\beta,\ldots,\beta),\qquad
 \beta=\sqrt{\frac{n-4}{n-1}}.                           \label{eq:prec-state-pair}
\end{equation}
Both norms are $\sqrt n$, and
\begin{equation}
 \frac1{\sqrt n}\leq
 \norm{v^{(0)}/\sqrt n-v^{(1)}/\sqrt n}_2
 \leq\sqrt{2/n}.                                        \label{eq:prec-close-states}
\end{equation}
Choose unitary extensions $U_0,U_1$ related by the minimum planar rotation
between the prepared states.  Then
$\norm{U_0-U_1}_2=\Theta(n^{-1/2})$, so the hybrid argument requires
$\Omega(\sqrt n)$ oracle calls for constant-bias distinction.

On the other hand, $O(A)$ controlled calls to $W_v$ estimate
$\operatorname{Re}\bra{0^a,1}W_v\ket{0^a,1}$ to additive error $1/(16A)$.
Multiplication by the reported $\alpha$ contributes error at most $1/16$,
and block-encoding error contributes at most $1/8$.  The first diagonal entry
is therefore known within $1/4$, distinguishing one from two.  The composed
distinguisher uses $O(qA)$ original oracle calls, proving
\eqref{eq:prec-qA}.

Rattew and Rebentrost give an exact constant-query block encoding of
$\diag(v/r_v)$ from $U_v$ \cite{rattew2023-non-linear-transformations-of-quantum}.
As an encoding of $\diag(v)$ it has normalization $r_v=\sqrt n$ on
\eqref{eq:prec-state-pair}, attaining the product in order.  This is
worst-case tightness, not a per-vector lower bound for every structured
preparation circuit.
\end{proof}

With copy-only access, the pair in \eqref{eq:prec-state-pair} requires
$\Omega(n)$ copies: one-copy overlap is $1-\Theta(1/n)$, so only
$k=\Omega(n)$ copies have constant trace distance.  Thus inverse access in
\cref{def:state-preparation-access} changes the exponent but does not erase
the interface cost.

The same obstruction sharpens near a strictly complementary endpoint.  Let
$n=2m$, $m\geq5$, $0<\mu\leq1$, and define
\begin{equation}
 \beta=\sqrt{\frac{m-4}{m-1}},\quad
 v=(2,\beta\mathbf1_{m-1}),\quad
 r=(m-1)(1-\beta)=\frac3{1+\beta},\quad
 z=2(r,\mathbf1_{m-1}).                                  \label{eq:prec-central-pair}
\end{equation}
For the one-sparse LP
\begin{equation}
 A=[I_m\;0],\qquad b=\mathbf1_m,\qquad c=(z,\mathbf1_m), \label{eq:prec-central-lp}
\end{equation}
take the common $x=(\mathbf1_m,\mu\mathbf1_m)$ and
\begin{align}
 s^{(0)}&=(\mu\mathbf1_m,\mathbf1_m),&
 y^{(0)}&=z-\mu\mathbf1_m,\nonumber\\
 s^{(1)}&=(\mu v,\mathbf1_m),&
 y^{(1)}&=z-\mu v.                                      \label{eq:prec-central-triples}
\end{align}
Both triples are feasible.  Since $\norm v=\norm{\mathbf1_m}$ and
$z^T(v-\mathbf1_m)=0$, the corresponding slack norms and dual norms agree.

\begin{theorem}[Strict-complementarity state-access barrier]
\label{thm:prec-centrality-state}
In controlled state-preparation access to the slack and dual states in
\eqref{eq:prec-central-triples}, distinguishing
\begin{equation}
 \norm{Xs-\bar\mu\mathbf1}_2\leq\bar\mu/4
 \quad\text{from}\quad
 \norm{Xs-\bar\mu\mathbf1}_2\geq3\bar\mu/4              \label{eq:prec-centrality-promise}
\end{equation}
requires
\begin{equation}
 \boxed{\Omega(\sqrt n/\mu)}                             \label{eq:prec-centrality-lb}
\end{equation}
total controlled preparation/inverse calls.  With copy-only access the bound
is $\Omega(n/\mu^2)$.  The result holds although both triples approach the
same strictly complementary optimum.  The fixed LP data, the common $x$, all
vector norms, positivity, and the real phase conventions are free.
\end{theorem}

\begin{proof}
The zero triple is exactly centered with $\bar\mu_0=\mu$.  For the other,
$\bar\mu_1<\mu$ while the first complementarity product is $2\mu$, so
\begin{equation}
 \norm{Xs^{(1)}-\bar\mu_1\mathbf1}_2>\mu>\bar\mu_1.       \label{eq:prec-central-gap}
\end{equation}
The normalized slack-state distance is $\Theta(\mu/\sqrt n)$.  The duals
differ by the same unnormalized vector $\mu(v-\mathbf1_m)$ and have common
norm $\Omega(\sqrt m)$, so their normalized-state distance is also
$O(\mu/\sqrt n)$.  Choose close minimum-rotation extensions for both
oracles.  A hybrid over queries to either oracle gives
\eqref{eq:prec-centrality-lb}.  For independent copies, fidelity
multiplicativity gives the square of this dependence.  Finally both triples
converge to
$((\mathbf1_m,0),z,(0,\mathbf1_m))$, proving strict complementarity.
\end{proof}

The triples can also be generated by a sparse LP's objective-state oracle.
Put $d=v-\mathbf1_m$,
$z_0=2(r,\mathbf1_{m-1})$, and $z_\mu=z_0-(\mu/2)d$.  Then
\begin{equation}
 z_0^Td=0,\qquad \mathbf1^Td=1-r,\qquad 1\leq\norm d_2^2\leq2.
 \label{eq:prec-input-identities}
\end{equation}
\begin{theorem}[Input-generated state-only barrier]
\label{thm:prec-input-generated}
Use the common one-sparse $A=[I_m\;0]$, $b=\mathbf1_m$, and objectives
\begin{equation}
 c^{(0)}=(z_\mu,\mathbf1_m),\qquad
 c^{(1)}=(z_\mu+\mu d,\mathbf1_m).                       \label{eq:prec-input-objectives}
\end{equation}
Let the only input-dependent access be a controlled clean preparation of the
normalized objective and its inverse; the common norm, $A$, and $b$ are free.
Distinguishing the inputs, certifying which one is centered at the common pair
$x=(\mathbf1_m,\mu\mathbf1_m)$,
$y=z_\mu-\mu\mathbf1_m$, or estimating the optimum to additive error
$\mu/8$, needs $\Omega(\sqrt n/\mu)$ controlled oracle calls.  Copy-only
access needs $\Omega(n/\mu^2)$ copies.

Moreover, multiply both objectives symmetrically by $n$, take
$b=\mathbf1_m/6$, and use
\[
 \widetilde c^{(0)}=n(z_0-\tfrac\mu2d,\mathbf1_m),\qquad
 \widetilde c^{(1)}=n(z_0+\tfrac\mu2d,\mathbf1_m).
\]
The normalized objective oracles do not change, the value gap obeys
\begin{equation}
 |\operatorname{OPT}_1-\operatorname{OPT}_0|
 =\frac{n\mu}{6}(r-1)\geq\frac{n\mu}{12},                \label{eq:prec-value-gap}
\end{equation}
and additive value accuracy $n\mu/48$ has the same query lower bound.  This
is a fixed fraction of the centered total duality gap $n\mu$; here
$n=\nu$ is the canonical LP barrier parameter of \cref{sec:prelim-conic}.
\end{theorem}

\begin{proof}
The objective norms agree because
$2\mu z_\mu^Td+\mu^2\norm d^2=0$, and their normalized-state distance is
$\Theta(\mu/\sqrt n)$.  Minimum-rotation objective-state oracles therefore
require $\Omega(\sqrt n/\mu)$ calls to distinguish; fidelity
multiplicativity gives $\Omega(n/\mu^2)$ copies.  The displayed common pair
is centered for input zero and violates \eqref{eq:prec-centrality-promise}
for input one.  Every feasible point has first block $\mathbf1_m$ and its
positive-cost second block vanishes at optimum, so the original value gap is
$\mu|\mathbf1^Td|\geq\mu/2$.  Additive error $\mu/8$ distinguishes the
inputs.

For the scaled claim, the common point
\[
 \widetilde x=(\mathbf1_m/6,(\mu/n)\mathbf1_m),\qquad
 \widetilde y=n(z_0-\tfrac\mu2d)-6\mu\mathbf1_m
\]
has every complementarity product equal to $\mu$ for input zero; input one
remains positive because $n(1-\beta)\leq5<6$ and lies outside every fixed
$N_2$ neighborhood of radius below one.  Its centered total duality gap is
$n\mu$.  Feasibility fixes the first primal block to $\mathbf1_m/6$, which
gives \eqref{eq:prec-value-gap}; accuracy $n\mu/48$ distinguishes the inputs.
The inverse-$\mu$ factor is therefore not caused by demanding accuracy finer
than the centered duality-gap scale.
\end{proof}

This input-generated theorem has a sharp boundary.  A binary coefficient-value
oracle exposes the known exceptional coordinate, whose values differ by
$\Theta(\mu)$, in one query at sufficient requested precision.  Distinct
value-register strings are not close unitaries merely because the encoded
reals are close.  Likewise, the naturally normalized Newton diagonals for the
two inputs differ by a constant in their first entry, so a freely supplied
Newton block encoding distinguishes them in $O(1)$ calls.  The
$\Omega(\sqrt n/\mu)$ theorem is therefore an end-to-end result only for the
state-only objective/update interface; it does not lift to raw exact-value or
precompiled-Newton access.

State access also obstructs fraction-to-boundary line search.

\begin{corollary}[Fraction-to-boundary]
\label{cor:prec-fraction-boundary}
For $n\geq4$, let $x=\mathbf1_n$ and compare the equal-norm directions
\[
 \delta^{(0)}=-\tfrac12\mathbf1_n,\qquad
 \delta^{(1)}=(-1,-\gamma,\ldots,-\gamma),\qquad
 \gamma=\tfrac12\sqrt{\frac{n-4}{n-1}}.
\]
Their maximal positive steps are two and one, while their normalized states
are $\Theta(n^{-1/2})$ apart.  A state-only rule that is never larger than one
on the second input but is at least $3/2$ on the first needs
$\Omega(\sqrt n)$ controlled state-preparation calls.
\end{corollary}

\begin{proof}
Choose the two preparation unitaries as minimum rotations.  Any valid output
distinguishes them with constant bias, while a $Q$-query hybrid has distance
$O(Q/\sqrt n)$.  Thus $Q=\Omega(\sqrt n)$.
\end{proof}

The next theorem rules out bypassing diagonal compilation by directly solving
the following easy system.  For $n\geq4$, fix $0<\eta\leq1/4$ and define
\begin{equation}
 v^{(0)}=\mathbf1_n,\qquad
 v^{(1)}=(1+\eta,\beta_\eta,\ldots,\beta_\eta),\qquad
 \beta_\eta=\sqrt{1-\frac{2\eta+\eta^2}{n-1}}.           \label{eq:prec-transition-pair}
\end{equation}
Both norms equal $\sqrt n$, and their normalized distance $d_\eta$ satisfies
\begin{equation}
 \frac\eta{\sqrt n}\leq d_\eta\leq\frac{2\eta}{\sqrt n}. \label{eq:prec-transition-distance}
\end{equation}
Let $D_b=\diag(v^{(b)})$,
$r=(e_1+e_2)/\sqrt2$, and
$\ket{z_b}=D_b^{-1}r/\norm{D_b^{-1}r}_2$.  These systems are one-sparse and
$\kappa(D_b)<2$.

\begin{theorem}[State to next solve]
\label{thm:prec-state-next-solve}
Give an algorithm controlled access to the minimum-rotation preparations
$U_b,U_b^\dagger$ of the vectors in \eqref{eq:prec-transition-pair}.  If its
output $\rho_b$ satisfies
\begin{equation}
 \Dtr(\rho_b,\ket{z_b}\!\bra{z_b})\leq\eta/12            \label{eq:prec-next-solve-error}
\end{equation}
for both inputs, then it uses $\Omega(\sqrt n)$ oracle calls.  With copy-only
iterate access it needs $\Omega(n)$ copies.  In particular, fixed
$\eta=1/4$ is a constant-accuracy lower bound for a condition-below-two
system.
\end{theorem}

\begin{proof}
For $b=0$, measuring coordinate one in $\ket{z_b}$ has probability $1/2$.
For $b=1$, put $q=\beta_\eta/(1+\eta)$.  The probability is
$q^2/(1+q^2)$, and
\begin{equation}
 \frac12-\frac{q^2}{1+q^2}
 \geq\frac12-\frac1{1+(1+\eta)^2}\geq\frac\eta3.         \label{eq:prec-solution-separation}
\end{equation}
Contractivity and \eqref{eq:prec-next-solve-error} therefore make the two
algorithm outputs at least $\eta/6$ apart in trace distance.  A $Q$-query
hybrid is at most $Qd_\eta\leq2Q\eta/\sqrt n$ apart, whence
$Q\geq\sqrt n/12$.  For $k$ copies the input trace distance is at most
$\sqrt{k}d_\eta=O(\sqrt{k}\eta/\sqrt n)$, forcing $k=\Omega(n)$.
\end{proof}

The Hermitian dilation
$\left(\begin{smallmatrix}0&D_b\\D_b&0\end{smallmatrix}\right)$ has the
same sparsity and condition number, so this is also a standard Hermitian QLS
instance.  It does not contradict a polylogarithmic-dimensional QLS
algorithm: such an algorithm assumes the matrix oracle that this theorem
requires one to construct or bypass from the iterate state.

Per-round costs add only under an explicit fresh-oracle contract.  At round
$t=1,\ldots,T$, after earlier outputs are committed, an adversary chooses a
fresh bit $b_t$ independent of the current workspace and supplies only a new
minimum-rotation pair $U_{t,b_t},U_{t,b_t}^\dagger$ for
\eqref{eq:prec-transition-pair}; future oracles are unavailable.  Before the
next round, the algorithm must either satisfy
\eqref{eq:prec-next-solve-error}, or, after $S_t$ setup queries, expose a
reusable controlled $(\alpha_t,a,\eta/8)$ block encoding of $D_{b_t}$ with
$\alpha_t\leq A_t$.  One invocation may use another $q_t$ state-oracle calls.

\begin{theorem}[Online fresh-oracle direct sum]
\label{thm:prec-online-direct-sum}
In the preceding interface, every fixed worst-case per-round budget satisfies
\begin{equation}
 \sum_{t=1}^TQ_t=\Omega(T\sqrt n)                         \label{eq:prec-online-solves}
\end{equation}
for next-solution-state outputs, and
\begin{equation}
 \boxed{\sum_{t=1}^T(\eta S_t+q_tA_t)=\Omega(T\sqrt n)}  \label{eq:prec-online-compilers}
\end{equation}
for reusable diagonal compilers.  For fixed short-step $\eta$, the latter is
equivalently $\sum_t(S_t+q_tA_t)=\Omega(T\sqrt n)$.
\end{theorem}

\begin{proof}
Condition on the complete workspace and transcript at the start of round
$t$.  Freshness makes that side information independent of $b_t$, so it is a
fixed ancilla in \cref{thm:prec-state-next-solve}, which gives
$Q_t=\Omega(\sqrt n)$.  For a compiler, run setup once and estimate the first
encoded entry to $O(\eta/A_t)$ using $O(A_t/\eta)$ invocations.  This
distinguisher makes $S_t+O(q_tA_t/\eta)$ state-oracle calls, but the close
input oracles require $\Omega(\sqrt n/\eta)$ calls.  Hence
$\eta S_t+q_tA_t=\Omega(\sqrt n)$ in each conditioned round.  Summation
proves both claims.
\end{proof}

The full oracle ledger in \cref{thm:prec-online-direct-sum} is load-bearing:
\begin{itemize}[leftmargin=*]
 \item The atomic charged operations are controlled calls to the current
 $U_{t,b_t}$ and its inverse.  Gate complexity is not lower-bounded.
 \item The common norm $\sqrt n$, the two vector formulas, $r$, and the round
 number are free; only $b_t$ is hidden.
 \item No coordinate-value, QRAM, objective-coefficient, or compiled matrix
 oracle is available.  Each could reveal the exceptional coordinate in one
 query.
 \item The complete off-zero action is promised: the adversary may choose the
 close minimum-rotation completions.  Specifying only $U\ket0$ would not
 define a query lower bound.
 \item One-time setup and reusable advice are charged through $S_t$, while
 per-invocation use is charged through $q_t$.  Advice consumed by calls must
 be prepared in sufficient quantity and cannot be cloned for free.
 \item The required output is committed before the next fresh oracle arrives;
 this causality makes the conditional bounds additive.
\end{itemize}
A heralded compiler with success probability bounded below obeys the same
ledger after charging repetitions.  An unheralded constant failure
probability is not covered when $\eta\to0$, because it can exceed the target
states' $\Theta(\eta)$ separation.

The hard stream remains positive and obeys a coordinatewise short-update
promise: $|v_i^{(b_{t+1})}-v_i^{(b_t)}|\leq\eta v_i^{(b_t)}$.  Thus
positivity and short relative motion do not defeat the online theorem.  They
also do not turn it into a fixed-instance QIPM iteration bound.  The exact
obstructions are:
\begin{enumerate}[leftmargin=*]
 \item \emph{Freshness is stronger than a fixed optimization instance.}
 Every QIPM iterate is determined by one static input and prior computation;
 no independent oracle normally arrives after each output.
 \item \emph{One-vector worst-case hardness is not automatically a direct
 sum.}  If every iterate depends on one bit, learning it once makes every
 later diagonal easy.
 \item \emph{Circuit generation must be expanded to atomic queries.}  A
 circuit for $U_{x_{t+1}}$ may call $U_{x_t}$, which contains the prior
 history.  Unit-cost treatment hides inlining, while charging both as
 independent primitives double-counts it.
 \item \emph{Off-zero leakage is model dependent.}  The hybrid chooses close
 extensions, whereas a concrete solver circuit may reveal coordinate or
 trajectory data away from $\ket0$.  Its entire oracle action must be fixed.
 \item \emph{All-round coherent access changes the direct-sum problem.}  If a
 single multiplexed query reaches every time-indexed oracle, the causal proof
 fails and a separate coherent direct-sum theorem is needed.
 \item \emph{The QIPM equations may permit a bypass.}  The diagonal systems
 above have not been embedded as consecutive Newton systems of one fixed
 sparse LP; cancellation, scale-free formulations, or history-state methods
 might avoid separate compilation.
 \item \emph{The output contract matters.}  An architecture that exposes a
 reusable diagonal or a Newton state every round meets the theorem's
 contract.  An algorithm that proves its step analytically and emits only one
 final observable need not.
\end{enumerate}
Therefore a fixed-LP $\Omega(T\sqrt n)$ result would require $T$ independent
hard pieces in one sparse instance, a trajectory that reveals them one by
one, and a direct-sum theorem against the static input oracle.  None is proved
here.

Finally, coherent update errors accumulate.  If
\begin{equation}
 \norm{\ket{\widetilde v}-\ket{v/r_v}}_2\leq\xi,\qquad
 |\widetilde r_v-r_v|\leq\delta_r,                        \label{eq:prec-update-input-error}
\end{equation}
amplitude-to-diagonal lifting has absolute operator error at most
$\delta_r+\widetilde r_v\xi$.  For coherent LCU updates
$X_{t+1}=X_t+h_t\Delta_t$, let
$E_t:=\norm{\widetilde X_t-X_t}_2$ be the accumulated absolute operator error
before update $t$.  If the step lengths and LCU arithmetic are exact and the
only new error is the lifted direction error, the operator-norm triangle
inequality gives
\begin{equation}
 E_{t+1}\leq E_t+|h_t|
 (\delta_{r,t}+\widetilde r_t\xi_t),\qquad
 E_T\leq E_0+\sum_{t<T}|h_t|
 (\delta_{r,t}+\widetilde r_t\xi_t).                     \label{eq:prec-update-error-ledger}
\end{equation}
Rounding, imperfect LCU coefficients, or postselection errors must be added as
separate terms.
Late-path centrality therefore constrains the cumulative weighted error, not
only each local preparation.  Literal circuit reuse nests as well.  Let $C_t$
be the cost of the current iterate-preparation circuit, let $Q_t$ be the
number of its calls made by the next direction circuit, and let $G_t$ be all
nonrecursive work in that update.  With literal inlining, additive uniform
gate/query costs, and no sharing or memoization, the exact accounting is
\begin{equation}
 C_{t+1}=G_t+(1+Q_t)C_t.                                 \label{eq:prec-circuit-accounting}
\end{equation}
This identity is for literal composition, not a universal lower bound against
direct structured or history-state oracles.

\subsection{Lazy central oracles and the caching ceiling}
\label{sec:preconditioning-lazy}

Changing all diagonal weights does not imply that they must all be rewritten.
The following fixed sparse LP makes this failure explicit.  For $K$ disjoint
blocks of length $L$, let
$d_{a,i}=u_{a,i}-v_{a,i}$ and impose
\begin{equation}
 d_{a,0}=1,\qquad
 d_{a,i}-\varsigma_{a,i}d_{a,i-1}=0\quad(1\leq i\leq L),  \label{eq:prec-lazy-paths}
\end{equation}
with all $u,v\geq0$.  Minimize
\begin{equation}
 \sum_{a=1}^K\left[
  \sum_{i=0}^L(u_{a,i}+v_{a,i})
  +\alpha(u_{a,L}-v_{a,L})\right],\qquad 0<\alpha<1.       \label{eq:prec-lazy-objective}
\end{equation}
Rows and columns have constant sparsity and the instance size is
$\Theta(KL)$.  With
\begin{equation}
 p_{a,0}=1,\qquad p_{a,i}=\prod_{j=1}^i\varsigma_{a,j},\qquad
 q(\mu)=\mu+\sqrt{\mu^2+1},                              \label{eq:prec-lazy-prefix}
\end{equation}
the primal central path is exactly
\begin{equation}
 (u_{a,i}(\mu),v_{a,i}(\mu))
 =\left(\frac{q(\mu)+p_{a,i}}2,
         \frac{q(\mu)-p_{a,i}}2\right).                  \label{eq:prec-lazy-factorization}
\end{equation}
Indeed feasibility fixes the difference $p_{a,i}$; differentiating the
one-dimensional pair barrier gives the common sum $q(\mu)$.  In the fixed
orthonormal null basis
$(e_{u_{a,i}}+e_{v_{a,i}})/\sqrt2$, the reduced log-barrier Hessian
$W^\top X^{-2}W$ is the scalar
\[
 2[(q(\mu)+1)^{-2}+(q(\mu)-1)^{-2}]I,
\]
so it has condition one throughout the trajectory.  All changing numerical
information is the public scalar $q(\mu)$; all hidden structural information
is the static prefix table.

Give reversible raw access to bits
$\varsigma_{a,i}=(-1)^{b_{a,i}}$.  An endpoint central lookup at block $a$
returns \eqref{eq:prec-lazy-factorization} for $i=L$ and therefore reveals
$P_a=p_{a,L}$, even with coordinate error strictly below $1/2$.

\begin{theorem}[Lazy endpoint lookup]
\label{thm:prec-lazy-lookup}
One exact coherent endpoint lookup, or one whose decoded sign is correct with
probability at least $2/3$, needs $\Omega(L)$ raw queries in the worst case,
even when the block index is in superposition.  The bound is tight.
\end{theorem}

\begin{proof}
Fix the block register to any basis value $a$ and fix all other blocks.
Decoding the returned pair computes parity of the remaining $L$ bits, which
has bounded-error quantum query complexity $\Omega(L)$.  Conversely, the
exact quantum parity algorithm computes the selected block's parity in
$O(L)$ queries while carrying $a$ coherently as a spectator; reversible
arithmetic then combines it with $q(\mu)$.  The same circuit works in
superposition over $a$.
\end{proof}

Call a resource an exact reusable endpoint oracle if, once constructed, a
fixed processor can obtain every one of the $K$ endpoint values without
further raw queries, and the resource survives or otherwise supports those
$K$ invocations.

\begin{theorem}[Reusable endpoint preprocessing]
\label{thm:prec-reusable-endpoints}
Constructing an exact reusable endpoint oracle needs at least
\begin{equation}
 \boxed{\left\lceil\frac{KL}{2}\right\rceil}             \label{eq:prec-reusable-cost}
\end{equation}
raw queries.  A bounded-error resource also needs $\Omega(KL)$ queries when
all $K$ successive answers are jointly correct with probability at least
$2/3$ on every input.
\end{theorem}

\begin{proof}
After preprocessing, invoke the resource on every block and multiply the
decoded signs.  The result is
\[
 \prod_{a=1}^KP_a=(-1)^{\bigoplus_{a=1}^K\bigoplus_{i=1}^Lb_{a,i}},
\]
the parity of all $KL$ raw bits.  Exact quantum parity requires
$\lceil KL/2\rceil$ queries, and bounded-error parity requires $\Omega(KL)$.
The proof permits entangled program states and implicit formulas; reusability,
not classical table materialization, is the decisive promise.
\end{proof}

\begin{theorem}[Explicit-coordinate trajectory service]
\label{thm:prec-lazy-trajectory}
Fix arbitrary parameters $\mu_1,\ldots,\mu_K>0$.  A service that at stage
$a$ returns the classical endpoint pair of block $a$ at $\mu_a$, with both
coordinate errors below $1/2$ and all answers jointly correct with probability
at least $2/3$ on every input, uses
\begin{equation}
 \boxed{\Omega(KL)}                                      \label{eq:prec-trajectory-cost}
\end{equation}
total raw queries across preprocessing and all stages.  Thus its amortized
cost is $\Omega(L)$.
\end{theorem}

\begin{proof}
Each returned pair reveals the independent $P_a$.  Multiplying the $K$
answers computes parity of all $KL$ input bits, so bounded-error quantum
parity proves \eqref{eq:prec-trajectory-cost}.
\end{proof}

Theorem~\ref{thm:prec-lazy-trajectory} is a lower bound for an explicit-cell
service contract, not for every QIPM.  Quantum cell-probe lower bounds likewise
require a specified storage-and-probe model and do not transfer automatically
to an analytic central path \cite{sen2001-lower-bounds-quantum-cell-probe}.
There is no unconditional per-IPM-iteration refresh theorem here, for five
separate reasons:
\begin{enumerate}[leftmargin=*]
 \item A fixed LP has no fresh per-iteration information.  Once all $KL$
 bits have been processed, unrestricted persistent memory can amortize that
 setup over arbitrarily many iterations.
 \item The path parameter may be low dimensional.  In
 \eqref{eq:prec-lazy-factorization} every changing coordinate is generated by
 one scalar $q(\mu)$, so a stored structural oracle needs only $O(1)$ extra
 arithmetic when $\mu$ changes.
 \item Cell-probe lower bounds need a query contract.  An IPM follows an
 analytic, algorithm-chosen trajectory rather than an adversarial sequence of
 cell updates and requests; \cref{thm:prec-lazy-trajectory} applies only
 after the distinct endpoint outputs are required.
 \item Coherent lookup defeats a naive block direct sum.  One $O(L)$ circuit
 evaluates the selected block in superposition.  A factor $K$ requires a
 classical-output contract or a suitable direct-product theorem.
 \item A normalized state is weaker than a data structure.  It need not
 support reusable random access to every coordinate; assuming that it does
 silently adds tomography or QRAM construction.
\end{enumerate}
This is the same caching ceiling as \eqref{eq:total-query-ledger} and the
static ledger following \cref{thm:lp-gain-plateau}: caching all signs can make
later raw-query cost vanish.  A genuine refresh lower bound would require a
memory cap, mandatory coordinate monitoring, nonshareable independent
invocations, or a proven online-hard fixed-LP trajectory.  Sparsity and
changing diagonal weights alone provide none of these.

\subsection{Synthesis}
\label{sec:preconditioning-synthesis}

Across all five interfaces, the invariant is a product or a ledger.  Separate
exact factor access gives $\alpha_P\sqrt K=\Omega(N)$ and, with reusable
compilation, $Ns+q\alpha=\Omega(N^2)$.  Scalar congruence obeys
$\sqrt\kappa\,\chi\geq c/\mu$; equality in order holds exactly for the
maximal class $r_\mu=\Omega(\mu)$ and $r_\mu=O(1)$.  The parity instance
preserves the total setup--RHS--solve--recovery cost even when a perfect factor
makes one phase trivial.  State preparation obeys $qA=\Omega(\sqrt n)$, and its iterative
version becomes additive only under a fresh-oracle contract.  Lazy evaluation
then shows why static preprocessing can defeat any unqualified per-iteration
claim.

These are access-model obstructions, not a verdict against preconditioning or
state reuse.  The positive modules in \cref{sec:positive-modules} are designed
to respect the same ledgers: they directly assemble favorable operators,
identify which projectors or cached structures are charged, and state what
output interface is actually produced.
